\documentclass{article}

\usepackage[preprint,nonatbib]{neurips_2020}

\usepackage[utf8]{inputenc}
\usepackage[T1]{fontenc}
\usepackage{hyperref}
\usepackage{url}
\usepackage{booktabs}
\usepackage{amsfonts}
\usepackage{microtype}
\usepackage{amsmath}
\usepackage{amssymb}

\title{Anticipatory Reinforcement Learning for Long-Lived Robot Planning}

\author{%
  Raghav Arora \\
  Department of Electrical and Computer Engineering \\
  The University of Texas at Austin \\
  \texttt{raghavaurora@utexas.edu}
}

\begin{document}

\maketitle

\begin{abstract}
Long-lived service robots act in persistent environments where solving the current task can change the cost of future tasks. Prior anticipatory-planning work uses offline learned estimates of expected future cost to bias symbolic planning, but largely remains planner-centric and short-horizon. We instead formulate anticipatory planning as a continual-task Markov Decision Process in which preparedness becomes the value of the persistent state under future tasks. We evaluate anticipatory RL against myopic RL on a symbolic restaurant benchmark aligned with prior restaurant settings and on a 5$\times$5 gridworld diagnostic. Current results show clear gains in the restaurant benchmark but mixed outcomes in the gridworld ablation.
\end{abstract}

\section{Introduction}
Robots deployed in homes, offices, and restaurants receive long sequences of instructions in environments that persist between tasks. In this setting, a myopic policy can create side effects: the cheapest way to finish the current task may move objects into states that make likely follow-up tasks harder, or into a \textit{trap} state. Anticipatory behavior instead seeks states that complete the present task while leaving the environment well prepared for what comes next.

Dhakal et al.\ introduced anticipatory planning as the joint optimization of current-task cost and expected future-task cost, with a learned offline estimator used to guide symbolic planning \cite{dhakal2023anticipatory}. Follow-up work extended this idea to task-and-motion planning \cite{dhakal2024tamp} and to larger home-like and restaurant settings \cite{talukder2024large}, while LLM-based systems anticipate future tasks before invoking a classical planner \cite{arora2024anticipate,singh2024anticipate}. These approaches are effective, but they assume a fixed task distribution and typically optimize only for the next task. Our thesis is that the same problem can be written directly as continual reinforcement learning over an augmented state $(s,\tau)$, letting preparedness emerge from value propagation across a persistent task stream.

\section{Problem Statement}
Dhakal et al.\ define anticipatory planning for a current task $\tau$ as choosing a goal-reaching end state that trades off immediate planning cost and expected cost of a sampled follow-up task:
\begin{equation}
\label{eq:dhakal}
s_g^*=\operatorname*{argmin}_{s_g' \in \mathcal{G}_\tau}\left[C^*_{s_g'}(s_0)+\sum_{\tau'} P(\tau')C^*_{\tau'}(s_g')\right].
\end{equation}
Here, $s_0$ is the current state, $\tau$ the current task, $\mathcal{G}_\tau$ the set of goal states satisfying $\tau$, $s_g'$ a candidate end state, $C^*_{s_g'}(s_0)$ the optimal cost of reaching $s_g'$ from $s_0$, $\tau'$ a follow-up task, $P(\tau')$ its probability, and $C^*_{\tau'}(s_g')$ the optimal cost of completing $\tau'$ from $s_g'$. Thus $s_g^*$ is the task-completing state that best balances immediate and expected future cost.\footnote{We write cost as $C$ rather than $V$ to avoid confusion with the RL value function below.} Prior work learns the anticipatory term \( C_{A.P}(s_g') = \sum_{\tau'} P(\tau')C^*_{\tau'}(s_g') \) offline with a GNN and uses it as a symbolic-planning heuristic \cite{dhakal2023anticipatory}.

We instead model the long-lived robot as an infinite-horizon MDP. Let $s \in \mathcal{S}$ be the physical state, $a \in \mathcal{A}$ the action, and $\tau \in \mathcal{T}$ the current task sampled from a task process $P(\tau)$. With goal predicate $g(s,\tau)\in\{0,1\}$, the augmented state is $x=(s,\tau)\in\mathcal{X}=\mathcal{S}\times\mathcal{T}$ and the task changes only at task boundaries:
\begin{equation}
\label{eq:task-transition}
\tau'=
\begin{cases}
\tau, & g(s',\tau)=0,\\
\text{sampled from } P(\tau), & g(s',\tau)=1.
\end{cases}
\end{equation}
The task-conditioned value function is
\begin{equation}
\label{eq:value}
V^\pi(s,\tau)=\mathbb{E}^\pi\!\left[\sum_{t=0}^{\infty}\gamma^t r(s_t,\tau_t,a_t,s_{t+1}) \mid s_0=s,\tau_0=\tau\right],
\end{equation}
and the anticipatory value of a physical state is
\begin{equation}
\label{eq:ap-value}
\mathbf{V}_{\mathrm{AP}}^\pi(s)=\mathbb{E}_{\tau \sim P(\tau)}[V^\pi(s,\tau)].
\end{equation}
Here, $\gamma$ is the discount factor controlling the effective anticipation horizon. Under this view, preparedness is the expected value of the persistent state under future tasks.
We define two agents, an anticipatory and a myopic agent. Both these agents use the same environment, action space, reward function, task generator, and reset schedule. They differ only in the Bellman target used at successful task completion:
\begin{equation}
\label{eq:targets}
y_{\text{ant}}=r+\gamma \max_{a'}Q(s',\tau',a'),
\qquad
y_{\text{myopic}}=r.
\end{equation}
Thus the anticipatory agent propagates value through the next sampled task, whereas the myopic agent treats success as terminal for learning while still acting in the same persistent world state distribution.

The benchmarks are procedurally generated task streams rather than an offline corpus. The first is a symbolic restaurant benchmark with recurring service and cleanup tasks (\texttt{serve\_water}, \texttt{make\_coffee}, \texttt{serve\_fruit\_bowl}, \texttt{clear\_containers}, \texttt{wash\_objects}), aligned with the restaurant task families used by Talukder et al.\ \cite{talukder2024large}. The second is a 5$\times$5 gridworld with four objects, three receptacles, and move/clear tasks. We evaluate greedy matched-seed rollouts using success rate, average task steps, and average task return.

\section{Literature Review}
Anticipatory planning was introduced by Dhakal et al.\ as a planner-centric objective with an offline GNN estimator of follow-up-task cost \cite{dhakal2023anticipatory}. Dhakal et al.\ later extended the idea to anticipatory task and motion planning \cite{dhakal2024tamp}, while Talukder et al.\ scaled it to large home-like and restaurant settings \cite{talukder2024large}. In parallel, Arora et al.\ and Singh et al.\ use language models to anticipate future tasks before classical planning \cite{arora2024anticipate,singh2024anticipate}. On the RL side, UVFAs \cite{schaul2015uvfa}, successor features \cite{barreto2017successor}, future-task-aware side-effect avoidance \cite{krakovna2020avoiding}, and continual-RL surveys \cite{khetarpal2022continual} motivate transfer and persistence across tasks. Our formulation differs by treating anticipatory planning itself as value propagation across successful task boundaries.

\section{Technical Approach}
Our current implementation uses task-conditioned Deep Q-Network agents operating on the augmented state $(s,\tau)$ \cite{mnih2015dqn}. The restaurant benchmark uses a symbolic vector observation containing robot location, held object, object locations, contents, cleanliness, and task descriptor, with a discrete macro-action space for pick, place, wash, fill, brew, and fruit-serving actions. Its Q-function is approximated by a two-hidden-layer multilayer perceptron. The gridworld benchmark uses an image-like state encoding of the 2D grid plus the active task, with a six-action discrete space consisting of four moves, pick, and place. Its Q-function is approximated by a convolutional encoder followed by a two-layer MLP head. In both cases, the environment persists across task completion, and a new task is sampled while the physical state is retained, so the agent must solve the present instruction while shaping the next state distribution it will inherit.

The myopic baseline is the same persistent environment with the same reset schedule as anticipatory, but successful transitions are stored as terminal for learning. Once task-boundary bootstrapping is enabled, the discount factor $\gamma$ controls the effective anticipatory horizon: smaller values emphasize near-term tasks, while larger values more strongly value preparedness for the continuing stream.

Moving forward in the project, we will compare against supervised anticipatory-planning baselines in the style of Dhakal et al. We will evaluate RL against these supervised estimators under matched task sequences, first in the symbolic restaurant setting and then in a larger ProcTHOR-backed embodied simulation.

The full evaluation protocol will separate effect size from statistical significance. We will run multiple matched seeds, report confidence intervals for return and preparedness metrics, and use paired tests on per-sequence outcomes: paired bootstrap confidence intervals and permutation tests for continuous metrics, and McNemar-style tests for binary success outcomes when appropriate. The current milestone table is therefore descriptive evidence from saved runs, not a finalized statistical claim.

\section{Intermediate / Preliminary Results}
\begin{table}[t]
\centering
\small
\setlength{\tabcolsep}{4pt}
\begin{tabular}{lcccc}
\toprule
Metric & Rest.-Ant. & Rest.-Myo. & Grid-Ant. & Grid-Myo. \\
\midrule
Success rate (\%)   & 99.68 & 99.60 & 86.0  & 86.7 \\
Avg.\ task steps    & 1.79  & 1.94  & 33.1  & 31.8 \\
Avg.\ task return   & 11.46 & 10.79 & -38.53 & -37.01 \\
Reward / step       & 6.39  & 5.55  & -1.166 & -1.165 \\
\bottomrule
\end{tabular}
\caption{Current matched-seed greedy comparison runs. Restaurant metrics come from the saved symbolic restaurant comparison; gridworld metrics come from the current move-heavy grid diagnostic.}
\label{tab:results}
\end{table}

Table \ref{tab:results} highlights the results in both environments. In the restaurant benchmark, the anticipatory agent slightly improves success while reducing average task steps and substantially increasing reward per step, suggesting that cross-task bootstrapping helps when follow-up tasks are strongly coupled. In the gridworld setup, the anticipatory agent underperforms the myopic baseline on the main task metrics, indicating that the current implementation does not yet convert preparedness into consistently better overall performance. At this milestone, the continual-task MDP formulation appears promising in structured persistent settings, but optimization stability and sample efficiency remain open issues.

\clearpage
\phantomsection
\label{sec:refs}
\bibliographystyle{unsrt}
\bibliography{milestone_refs}

\end{document}
